% The plaintext guarantees; execution attacks are treated on the next slide.
\begin{frame}{Distillation: Observations}
  \centering
  \setbeamercolor{itemize item}{fg=CVQCTeal}
  \begin{tikzpicture}
    \path[use as bounding box] (-6.6,-3.15) rectangle (6.6,3.05);

    \node[anchor=north west,inner sep=0pt,text=CVQCInk] at (-6.35,2.45) {%
      \begin{minipage}{6.8cm}
        \raggedright
        \begin{itemize}
          \setlength{\itemsep}{0.75em}
          \item<1-> \textbf{Correctness:}\\[0.15em]
            $\ket{0}^{\otimes 4}\;\longmapsto\;\ket{0}$.
          \item<2-> \textbf{Soundness:} one wrong input\\[0.15em]
            $\longmapsto\;\ket{0}$ or reject.
          \item<3-> Can be nested recursively:\\[0.15em]
            $\ket{0^{4^L}}\;\longmapsto\;\ket{0}$.
          \item<4-> Generalizes to EPR pairs.
        \end{itemize}
      \end{minipage}%
    };

    \node[anchor=north east,inner sep=0pt,text=CVQCInk] at (6.35,2.65) {%
      \resizebox{5.2cm}{!}{%
        \let\tiny\scriptstyle
        \hspace*{0.55cm}% qcircuit does not measure stick labels.
        $\Qcircuit @C=1.4em @R=1.2em {
          \lstick{\ket{0}} & \qw & \ctrl{1} & \targ & \qw & \rstick{\ket{0}}\qw \\
          \lstick{\ket{0}} & \qw & \targ & \qw & \meter & \cw \\
          \lstick{\ket{0}} & \qw & \ctrl{1} & \ctrl{-2} & \meterB{X} & \cw \\
          \lstick{\ket{0}} & \qw & \targ & \qw & \meter & \cw
        }$\hspace*{0.55cm}%
      }%
    };
  \end{tikzpicture}

  \note<1>{
    \footnotesize
    First consider honest execution on plaintext states. C-dist denotes
    the three CNOTs; the full routine includes measurements and checks.
    Four zero inputs remain four zeros after these CNOTs. The Z outcomes
    on wires two and four must both be zero. The X outcome on wire three
    is random, and either value is accepted. Thus honest execution accepts
    with certainty and leaves zero on the first wire. The small diagram
    depicts this honest case throughout the slide.
  }
  \note<2>{
    \footnotesize
    Assume at most one input is wrong and the other three are ideal zeros.
    The exceptional input may be an arbitrary state, including one
    entangled with an external reference. A computational-basis one on
    any one input propagates to at least one checked Z outcome, so that
    component is rejected. Acceptance projects out the erroneous
    component and leaves zero on the surviving wire. No guarantee on
    arbitrary four-input garbage is claimed here, and these plaintext
    guarantees assume that the circuit and checks are performed honestly.
    Source: Technical Overview, basic 4-to-1 example.
  }
  \note<3>{
    \footnotesize
    Group inputs in fours and feed each surviving output into the next
    level. A depth-L four-ary tree combines $4^L$ inputs into one output.
    The displayed map is honest correctness on ideal zero inputs, not
    unconditional soundness for arbitrary noisy inputs. For the error
    analysis, every bad parent needs at least two bad children: a single
    bad child is removed or the block rejects. Therefore a bad root
    requires at least $2^L$ bad leaves. The paper combines this counting
    rule with the preparation stage's joint bad-support bound.
  }
  \note<4>{
    \footnotesize
    State the extension without proof. The paper also gives 4-to-1
    routines for EPR pairs and for Hadamard magic states. For EPR pairs,
    apply the same C-dist pattern to each half and check the appropriate
    stabilizer parities. Their correctness and one-bad-input properties
    are proved in the paper. The recursive construction applies to those
    routines as well. CCZ uses a separate distillation routine, so the
    four-to-one picture here is not a claim about CCZ distillation.
    Source: Distillation section, 4-to-1 circuits and measurement checks.
  }
\end{frame}
